\chapter{Introduction}
\label{Section:Introduction}

This thesis studies constrained optimization and online resource control
under different oracle and feedback models. The restrictions arise from
costly domain projections, resource availability that depends on past
actions, and service rates that must be learned from binary feedback.
In the optimization problems,
linear optimization constructs points in the original domain while
subgradient updates
use an auxiliary domain. In energy allocation and scalar storage control,
parameterizing allocations or policies permits first-order learning, while
battery feedback and action clipping enforce resource availability.
The guarantees bound objective gap,
constraint violation, or regret; the scalar storage guarantee allows
adversarial arrival and cost sequences. For transmission-rate adaptation,
confidence bounds based on binary feedback yield a bound on time-average
expected queue length.

\section{Models and performance criteria}


Table~\ref{tab:intro:model-comparison} compares the models, available
information, performance criteria, and feasibility requirements.


\begin{table}[htbp]
\centering
\caption[Models, feedback, and performance criteria]{Principal distinctions
between the research models. Formal assumptions are stated in each chapter.}
\label{tab:intro:model-comparison}
\begin{singlespace}
\small
\begin{tabular}{@{}>{\raggedright\arraybackslash}p{0.11\textwidth}>{\raggedright\arraybackslash}p{0.43\textwidth}>{\raggedright\arraybackslash}p{0.39\textwidth}@{}}
\toprule
Chapter & Model and information & Comparison and feasibility \\
\midrule
\ref{Section:projfree-domain} &
Convex objective; linear oracle on $\X$, subgradients on
$\Y\supseteq\X$, and easy projection onto $\Y$. &
Fixed minimizer over $\X$; output in $\X$, with auxiliary queries
permitted outside $\X$. \\[0.5em]
\ref{Section:projfree} &
Additional convex inequality constraints; linear oracle on $\X$;
objective and constraint subgradients and constraint values on $\Y$;
easy projection onto $\Y$. &
Upper bound on objective gap at $\bar x\in\X$; upper bound on constraint
violation if an optimal Lagrange multiplier exists. \\[0.5em]
\ref{Section:ChapEnergyharvest} &
Finite battery capacity; i.i.d.\ energy arrivals from an unknown
distribution, observed before allocation; the full convex allocation-cost
function is revealed afterward. &
Expected regret against a fixed allocation whose per-slot expenditure is at
most the mean energy arrival; only the controller must respect the available
battery energy on every slot. \\[0.5em]
\ref{Section:ChapControl} &
Leaky storage with retention coefficient $0<\alpha<1$ and enough capacity to
avoid saturation; arbitrary arrivals in $[0,1]$ and entire jointly
convex state--action cost functions revealed after acting. &
Regret against fixed policies that allocate fractions of past arrivals over
time; controller withdrawals never exceed the available resource. \\[0.5em]
\ref{Section:ChapRateAdaptation} &
Independent i.i.d.\ arrival and capacity sequences with unknown distributions;
known lower bound $\varepsilon$ on service slack. Rate in $[0,1]$ chosen
before observing capacity; binary acknowledgment even for an empty queue. &
A bound on time-average expected queue length that is independent of $T$. \\
\bottomrule
\end{tabular}
\end{singlespace}
\end{table}

\section{Contributions}

In Chapters~\ref{Section:projfree-domain} and~\ref{Section:projfree}, $T$
denotes the query budget for each oracle, including initialization.
In the online problems, $T$ denotes the number of rounds.

\subsection{Projection-free nonsmooth optimization}

Chapter~\ref{Section:projfree-domain} develops a method that separates the
construction of feasible points from the evaluation of objective
subgradients. Linear optimization generates points in the original convex
domain, while projected subgradient updates operate on an auxiliary domain
containing it. A multiplier couples the updates. The method returns a
feasible output with objective gap $O(T^{-1/2})$. Each update requires one linear
optimization call, one subgradient call, and one auxiliary projection.
The resulting oracle complexity is $O(\epsilon^{-2})$ for an objective
gap of at most $\epsilon$.

\subsection{Projection-free optimization with functional constraints}

Chapter~\ref{Section:projfree} formulates a joint bound on objective gap and
constraint violation. The output belongs to the original domain, and its
objective gap has an $O(T^{-1/2})$ upper bound. If an optimal Lagrange
multiplier exists, the constraint violation norm also satisfies an
$O(T^{-1/2})$ bound. The multiplier affects the guarantee but is not required
to choose the algorithm parameters.

\subsection{Online allocation with stochastic energy arrivals}

Chapter~\ref{Section:ChapEnergyharvest} develops an energy-allocation
controller by writing each allocation as the product of total energy usage
and a vector of allocation proportions. This representation permits
separate updates: a projected first-order update with battery feedback
controls total usage, while an entropic mirror-descent update controls the
proportions. The controller respects available energy on every slot and,
under the conditional-mean assumption~\eqref{eq:energy:conditional-mean},
achieves $O(\sqrt{T(1+\log n)})$ expected regret with sufficient battery
capacity $O(\sqrt T)$. The benchmark is the best fixed nonnegative allocation
whose total usage does not exceed the mean energy arrival. With the componentwise
subgradient bound fixed, the dependence on the number of subsystems
improves from $O(\sqrt n)$ in Yu and Neely~\cite{Yu2019} to
$O(\sqrt{\log n})$ for $n\ge2$, while the sufficient capacity is independent
of $n$.

\subsection{Online control of scalar storage}

Chapter~\ref{Section:ChapControl} introduces simplex disturbance-action
control (SDAC), which parameterizes storage policies by fractions of past
arrivals assigned to future withdrawals. Entropic mirror descent updates
the policy parameters, and action clipping limits each withdrawal to the
resource currently available.
For arbitrary arrivals in $[0,1]$ and full cost
feedback, with retention and cost parameters fixed, the controller achieves
$O(\sqrt{T\log(H+1)})$ regret against the
best fixed SDAC policy with memory $H$.

Choosing the memory length as a function of the horizon and bounding the
cost of truncation gives $O(\sqrt{T\log\log T})$ regret against the best
fixed infinite-memory SDAC policy
(Corollary~\ref{ch4:cor:regret-vs-infinite-sdac}). This policy class includes
every feasible scalar fixed-fraction policy under the same dynamics
(Corollary~\ref{ch4:cor:fraction-as-sdac}), so the guarantee also applies to
that comparison class.

For retention $\alpha\in(0,1)$ and positive state and action cost
Lipschitz bounds $L_x,L_u$, define
$G_\alpha\coloneqq\alpha(L_x/(1-\alpha)+L_u)$.
For every integer $H\ge1$ and horizon satisfying
\eqref{ch4:eq:sdac-lower-horizon}, and for every causal feasible controller,
including randomized controllers,
Theorem~\ref{ch4:thm:sdac-lower-bound} supplies a deterministic
arrival and cost sequence fixed before the interaction on which expected
regret against the fixed memory-$H$ SDAC class is at least
$G_\alpha\sqrt T/8$. The expectation is over the controller's
randomization. For fixed $H,\alpha,L_x,L_u$, the upper and lower bounds
therefore match in their $\sqrt T$ dependence on the horizon $T$.
Gaps remain in their dependence on memory length $H$ and retention
coefficient $\alpha$, as quantified in
Section~\ref{ch4:sec:regret-lower-bound}.

\subsection{Rate adaptation from binary feedback}

Chapter~\ref{Section:ChapRateAdaptation} develops and analyzes a
transmission-rate controller that applies UCB1 to a finite grid of rates
using binary success or failure feedback. A known lower bound
$\varepsilon>0$ on service slack determines the grid resolution. For
independent i.i.d.\ arrival and capacity sequences and an initially empty
queue, the controller achieves time-average expected queue length
$O((1+\log(1/\varepsilon))/\varepsilon^2)$ uniformly over all horizons. A
matching $\Omega(\varepsilon^{-2}\log(1/\varepsilon))$ lower bound on a suitable
deterministic instance and horizon shows that the logarithmic factor is
necessary for this UCB1 controller.

\section{Notation and conventions}
\label{sec:intro:notation}
\label{sec:notation}

We write $\Real_+$ for the nonnegative reals, $[n]\coloneqq\{1,\ldots,n\}$,
$\mathbf1$ for the all-ones vector, and $[v]_+$ for the componentwise positive
part. Vector inequalities are componentwise. All logarithms are natural unless
a base is shown. Empty sums are zero.
For an event $E$, its probability is $\mathbb P(E)$, and its indicator is $\ind{E}$.
We use $O(\cdot)$ for asymptotic upper bounds, $\Omega(\cdot)$ for
asymptotic lower bounds, and $\Theta(\cdot)$ for matching upper and lower
orders. The notation $a\coloneqq b$ defines $a$ as $b$.

Scalar variables and scalar coordinates are italic. In the abstract
optimization and geometry of Chapters~\ref{Section:Preliminaries}--\ref{Section:projfree},
plain italic letters also denote vectors, matrices, and linear maps; their
domains specify their type. In the resource models of
Chapters~\ref{Section:ChapEnergyharvest} and~\ref{Section:ChapControl}, bold
letters distinguish whole vectors and matrices from their scalar coordinates,
as in $\mathbf m_t=(m_{t,1},\ldots,m_{t,H})^\top$.
Randomness is
specified by the probability model and conditioning, rather than by case or
boldface alone.

Coordinate lists for column vectors are written $(v_1,\ldots,v_n)^\top$.
Scalar coordinates use subscripts: $v_i$ for coordinate $i\in[n]$ and
$v_{t,i}$ for that coordinate at round $t$. When time and coordinate indices
occur together, time comes first.
For vectors, $\inner{x}{y}\coloneqq x^\top y$.
An unqualified norm in a Euclidean argument is the Euclidean norm.

The unqualified symbol $\Pi_Q$ denotes Euclidean projection onto $Q$.

We write $\Reg_T(\mathcal C)$ for the controller's cumulative cost minus
the infimum comparator cost over $\mathcal C$, evaluated in hindsight on
the same realized inputs. Expected regret is $\E[\Reg_T(\mathcal C)]$.

\section{Organization of the thesis}

Chapter~\ref{Section:Preliminaries} collects the convex-analysis,
Bregman-divergence, and mirror-descent tools used in the research chapters
outlined above. The final chapter brings the results together and
identifies directions for further work.
